% GENERATED FILE. Edit canonical lessons, then run npm run generate:books.
% canonical-source-hash: fnv1a64:6579ece3be7c71a6
% canonical-lessons: GU-C122-rich, GU-C122-siyal, GU-C122-varu, GU-C122-kabutar, GU-C122-batak

\chapter{Bear, Jackal, Wolf, Pigeon, Duck}
\label{ch:gu-bear-jackal-wolf-pigeon-duck}
\hlchaptermodality{122}

\begin{chapteropening}
\textbf{By the end of this chapter:} \emph{I can say a bear, a jackal, a wolf, a pigeon and a duck.}

Complete the last lesson of chapter 122: i can say a bear, a jackal, a wolf, a pigeon and a duck.
\end{chapteropening}

\section[rī̃ch]{\gu{રીંછ} (rī̃ch) — a bear}

\label{lesson:GU-C122-rich}



\begin{quote}
Before the new one: say the Gujarati for a rabbit, then the Gujarati for a deer.
\end{quote}

\subsection*{You'll want to know: \gu{રીંછ}}
\textbf{\gu{રીંછ}} (\emph{rī̃ch}) — \textquotedblleft{}a bear\textquotedblright{}.

\begin{grammarlens}[title={What you've built}]
One more word for this chapter.
\end{grammarlens}

\subsection*{Your turn}
\begin{itemize}
\raggedright
  \item \emph{Say it:} \emph{rī̃ch}
  \item \emph{Say it:} \emph{rī̃ch}, once more
  \item \emph{Say it:} say \emph{rī̃ch}
  \item \emph{Recall:} say the Gujarati for a rabbit, then the Gujarati for a deer, then say \emph{rī̃ch} again
\end{itemize}

\begin{tcolorbox}[breakable,colback=teal!4,colframe=teal!35!black,title={Before you move on}]
What does \gu{રીંછ} mean? (\textquotedblleft{}A bear\textquotedblright{}.) Say it once more.
\end{tcolorbox}
\section[śiyāḷ]{\gu{શિયાળ} (śiyāḷ) — a jackal}

\label{lesson:GU-C122-siyal}



\begin{quote}
Before the new one: say the Gujarati for a deer, then the Gujarati for a bear.
\end{quote}

\subsection*{You'll want to know: \gu{શિયાળ}}
\textbf{\gu{શિયાળ}} (\emph{śiyāḷ}) — \textquotedblleft{}a jackal\textquotedblright{}.

\begin{grammarlens}[title={What you've built}]
One more word for this chapter.
\end{grammarlens}

\subsection*{Your turn}
\begin{itemize}
\raggedright
  \item \emph{Say it:} \emph{śiyāḷ}
  \item \emph{Say it:} \emph{śiyāḷ}, once more
  \item \emph{Say it:} say \emph{śiyāḷ}
  \item \emph{Recall:} say the Gujarati for a deer, then the Gujarati for a bear, then say \emph{śiyāḷ} again
\end{itemize}

\begin{tcolorbox}[breakable,colback=teal!4,colframe=teal!35!black,title={Before you move on}]
What does \gu{શિયાળ} mean? (\textquotedblleft{}A jackal\textquotedblright{}.) Say it once more.
\end{tcolorbox}
\section[varu]{\gu{વરુ} (varu) — a wolf}

\label{lesson:GU-C122-varu}



\begin{quote}
Before the new one: say the Gujarati for a bear, then the Gujarati for a jackal.
\end{quote}

\subsection*{You'll want to know: \gu{વરુ}}
\textbf{\gu{વરુ}} (\emph{varu}) — \textquotedblleft{}a wolf\textquotedblright{}.

\begin{grammarlens}[title={What you've built}]
One more word for this chapter.
\end{grammarlens}

\subsection*{Your turn}
\begin{itemize}
\raggedright
  \item \emph{Say it:} \emph{varu}
  \item \emph{Say it:} \emph{varu}, once more
  \item \emph{Say it:} say \emph{varu}
  \item \emph{Recall:} say the Gujarati for a bear, then the Gujarati for a jackal, then say \emph{varu} again
\end{itemize}

\begin{tcolorbox}[breakable,colback=teal!4,colframe=teal!35!black,title={Before you move on}]
What does \gu{વરુ} mean? (\textquotedblleft{}A wolf\textquotedblright{}.) Say it once more.
\end{tcolorbox}
\section[kabūtar]{\gu{કબૂતર} (kabūtar) — a pigeon}

\label{lesson:GU-C122-kabutar}



\begin{quote}
Before the new one: say the Gujarati for a jackal, then the Gujarati for a wolf.
\end{quote}

\subsection*{You'll want to know: \gu{કબૂતર}}
\textbf{\gu{કબૂતર}} (\emph{kabūtar}) — \textquotedblleft{}a pigeon\textquotedblright{}.

\begin{grammarlens}[title={What you've built}]
One more word for this chapter.
\end{grammarlens}

\subsection*{Your turn}
\begin{itemize}
\raggedright
  \item \emph{Say it:} \emph{kabūtar}
  \item \emph{Say it:} \emph{kabūtar}, once more
  \item \emph{Say it:} say \emph{kabūtar}
  \item \emph{Recall:} say the Gujarati for a jackal, then the Gujarati for a wolf, then say \emph{kabūtar} again
\end{itemize}

\begin{tcolorbox}[breakable,colback=teal!4,colframe=teal!35!black,title={Before you move on}]
What does \gu{કબૂતર} mean? (\textquotedblleft{}A pigeon\textquotedblright{}.) Say it once more.
\end{tcolorbox}
\section[batak]{\gu{બતક} (batak) — a duck}

\label{lesson:GU-C122-batak}



\begin{quote}
Before the new one: say the Gujarati for a wolf, then the Gujarati for a pigeon.
\end{quote}

\subsection*{You'll want to know: \gu{બતક}}
\textbf{\gu{બતક}} (\emph{batak}) — \textquotedblleft{}a duck\textquotedblright{}.

\begin{grammarlens}[title={What you've built}]
One more word for this chapter.
\end{grammarlens}

\subsection*{Your turn}
\begin{itemize}
\raggedright
  \item \emph{Say it:} \emph{batak}
  \item \emph{Say it:} \emph{batak}, once more
  \item \emph{Say it:} say \emph{batak}
  \item \emph{Recall:} say the Gujarati for a wolf, then the Gujarati for a pigeon, then say \emph{batak} again
\end{itemize}

\begin{tcolorbox}[breakable,colback=teal!4,colframe=teal!35!black,title={Before you move on}]
What does \gu{બતક} mean? (\textquotedblleft{}A duck\textquotedblright{}.) Say it once more.
\end{tcolorbox}
